% !TeX root = ../navier-stokes-manuscript.tex
% ============================================================
% 2.7 Compatibility and the Terminal Quantifier
% ============================================================

\subsection{Compatibility and the Terminal Quantifier}\label{sub:CompatibilityAndTheTerminalQuantifier}

Take the concrete rectangle \(\cR_{2,3}[u_0;T]\).
Its upper velocity row contains
\[
  V_0,V_1,V_2
  \in
  L^2(0,T;H^4),
\]
and its lower derivative row contains
\[
  V_1,V_2,V_3
  \in
  L^2(0,T;H^2).
\]
It also contains \(Q_0,Q_1,Q_2\) with their three differentiated momentum, divergence, and pressure equations.

Now retain temporal orders through \(n=1\) and lower the middle spatial height from \(H^3\) to \(H^2\).
The embeddings
\[
  H^4\hookrightarrow H^3,
  \qquad
  H^2\hookrightarrow H^1
\]
give
\[
  V_0,V_1
  \in
  L^2(0,T;H^3)
\]
and
\[
  V_1,V_2
  \in
  L^2(0,T;H^1).
\]
Retaining \(Q_0,Q_1\) and the corresponding equations gives exactly \(\cR_{1,2}[u_0;T]\).
Every surviving entry is the same derivative it was in the larger rectangle.
The operation changes the finite depth at which the history is viewed and preserves the history itself.

The general restriction has three independent cutoffs.
Let
\[
  N'\leq N,
  \qquad
  M'\leq M,
  \qquad
  0<T'\leq T<T_*.
\]
Define
\[
  \rho_{N',M';T'}^{N,M;T}:
  \cR_{N,M}[u_0;T]
  \longrightarrow
  \cR_{N',M'}[u_0;T']
\]
by retaining \(V_0,\ldots,V_{N'+1}\), retaining \(Q_0,\ldots,Q_{N'}\), using the continuous embeddings from the \(M\)-spaces to the \(M'\)-spaces, and restricting every retained function to \((0,T')\).
For every shared temporal and spatial index,
\[
  J_{n,\alpha}
  =
  \partial_t^n\partial_x^\alpha u,
  \qquad
  \Pi_{n,\alpha}
  =
  \partial_t^n\partial_x^\alpha p
\]
remain the same distributional derivatives.
The retained rows continue to satisfy the same pressure-complete mixed recurrence.

Restriction in stages agrees with direct restriction.
If
\[
  N''\leq N'\leq N,
  \qquad
  M''\leq M'\leq M,
  \qquad
  0<T''\leq T'\leq T<T_*,
\]
then
\[
  \rho_{N'',M'';T''}^{N',M';T'}
  \circ
  \rho_{N',M';T'}^{N,M;T}
  =
  \rho_{N'',M'';T''}^{N,M;T}.
\]
This is compatibility.
The same lower derivative is obtained from every larger rectangle that contains it, regardless of the intermediate cutoffs.
The coordinatewise intersections from the previous subsection refer to one velocity--pressure history.

Agreement among the rectangles is only half of what is needed at \(T_*\).
Fix \(T<T_*\).
Classical smoothness on the compact interval \([0,T]\) gives, for every finite \(n\) and \(q\),
\[
  \sup_{0\leq t\leq T}
  \norm{\partial_t^n u(t)}_{H^q}
  <
  \infty.
\]
Consequently,
\[
  \norm{\partial_t^n u}_{L^2(0,T;H^q)}^2
  \leq
  T
  \sup_{0\leq t\leq T}
  \norm{\partial_t^n u(t)}_{H^q}^2
  <
  \infty.
\]
Every finite rectangle is finite after its strict right endpoint has been fixed.
In quantifiers, this says
\[
  \forall\,0<T<T_*
  \ \forall\,N,M\in\Nzero
  \ \exists\,C_{T,N,M}<\infty:
  \qquad
  \mathfrak R_{N,M}(u;T)
  \leq
  C_{T,N,M}.
\]
The constant may grow as \(T\uparrow T_*\).

Whole-terminal control fixes the finite coordinate first and supplies one constant for every strict cutoff:
\[
  \forall\,N,M\in\Nzero
  \ \exists\,C_{N,M}(u_0,\nu,T_*)<\infty
  \ \forall\,0<T<T_*:
  \qquad
  \mathfrak R_{N,M}(u;T)
  \leq
  C_{N,M}(u_0,\nu,T_*).
\]
Equivalently,
\[
  \sup_{0<T<T_*}
  \mathfrak R_{N,M}(u;T)
  \leq
  C_{N,M}(u_0,\nu,T_*)
  <
  \infty
  \qquad
  \text{for every finite }(N,M).
\]
Each finite pair may have its own constant.
The uniformity concerns the approach of that fixed rectangle to \(T_*\).

The datum-generated theorem stated in \cref{thm:WholeTerminalCompatibleRectangleProducer} supplies this bound for every finite pressure-complete rectangle of the same history.
Section~3 takes their compatible family to a common endpoint.
